%1 協力者31名の選定
\begin{revisedblock}
We selected 31 experiment participants.
  %1-1 年齢帯
  The 31 experiment participants were in their 20s.
  %1-2 日常的なPC利用
  The 31 experiment participants used a PC daily.
\end{revisedblock}

%2トレーニングデータの取得方法
\deleted{LiDAR was used to acquire training data on the subject's posture while using the PC to perform four tasks.}
\revised{LiDAR was used to acquire training data on the experiment participant's posture while using the PC to perform four tasks.}\footnote{Our experiment was approved by the Ethics Committee of Shibaura Institute of Technology.}

\revised{Table~\ref{dataacquisition} summarizes data acquisition conditions for training data; each capture was repeated 20 times for mouse operation, pad operation, and typing, and 10 times for sitting still.}
\deleted{Table~\ref{trainacquisition} summarizes training data acquisition conditions; each capture was repeated 20 times for mouse operation, pad operation, and typing, and 10 times for sitting still.}
\begin{deletedblock}
As for mouse operation, the posture of the subjects was captured for 30 seconds as they manipulated a mouse to write the alphabet in order from ``A'' to ``Z''. The capturing was repeated 20 times to acquire training data. 
    Fig.~\ref{labels}\subref{pcmouse} shows the point cloud of the mouse operation. 
As for pad operation, the posture of the subjects was captured for 30 seconds as they performed a pad operation on the PC to write the alphabet in order from ``A'' to ``Z''. The capturing was repeated 20 times to acquire training data. 
    Fig.~\ref{labels}\subref{pcpad} shows the point cloud of the pad operation. 
As for sitting still, the posture of the subjects was captured for ten  seconds as they sat still without moving and did not touch the PC. The capturing was repeated ten times to acquire training data. 
    Fig.~\ref{labels}\subref{pcstay} shows the point cloud of sitting still. 
As for typing, the posture of the subjects was captured for 30 seconds as they typed the alphabet in order from ``A'' to ``Z''. The capturing was repeated 20 times to acquire training data. 
    Fig.~\ref{labels}\subref{pctyp} shows the point cloud of the typing. 
\end{deletedblock}

\begin{revisedblock}
\begin{table}[b]
  \centering
  \revised{\caption{Data acquisition conditions for each posture class.}}
  \deleted{\caption{Training data acquisition conditions for each posture class.}}
  \label{dataacquisition}
  \scalebox{0.9}{
  \begin{tabular}{l|l|c}
    \hline
    Posture class & Task & Duration \\ \hline
    Mouse operation & Write alphabet A--Z using mouse & 30 s \\
    Pad operation & Write alphabet A--Z using trackpad & 30 s \\
    Sitting still & Sit still without moving or touching PC & 10 s \\
    Typing & Type alphabet A--Z on keyboard & 30 s \\
    \hline
  \end{tabular}
  }
\end{table}
\end{revisedblock}

\revised{Fig.~\ref{labels} shows point clouds for each posture class.}

%6トレーニングデータの被験者
\revised{We randomly sampled three experiment participants from the 31 experiment participants.}
\deleted{The three subjects in this study are denoted as subjects A, B, and C.}
\revised{The three experiment participants in this study are denoted as experiment participants A, B, and C.}
    % Training data with subjects A, B, and C is called training data A, B, and C, and test data with subjects A, B, and C is called test data A, B, and C. 
    %6-1トレーニングデータの名称
    \deleted{Training data with subjects A, B, and C is called training data A, B, and C.}
    \revised{Training data with experiment participants A, B, and C is called training data A, B, and C.} 